    \section{Finite-Rank Shorting and Exact Fredholm Trace Cancellation}
    \label{sec:finite_rank}
    \label{sec:discussion}

We have resolved the rank-one block-operator obstruction: a compact
free covariance block need not be inverted, because the Schur subtraction is
the Anderson--Trapp shorted residual on the Cameron--Martin energy range.
Finite-rank targets require the same idea in operator-valued form, plus an
exact matrix/Fredholm identity for the Gaussian finite part.  Here Galerkin
approximation remains purely variational: truncation restricts the
Anderson--Trapp supremum to nested finite-dimensional test spaces.  The
limitation is infinite-rank observation, where Fredholm
determinants, trace-class perturbation hypotheses, and operator-valued finite
parts need separate treatment.

\subsection{Matrix algebraic extension}
\label{sec:finite_rank:matrix_extension}

For an $m$-dimensional observed subspace $E=\operatorname{span}\{\phi_1,\ldots,\phi_m\}$, the scalar target block $\Sigma_{\lambda,kk}$ is replaced by the covariance matrix $\Sigma_{\lambda,EE}\in\mathbb{R}^{m\times m}$. The scalar calibration becomes a matrix Lyapunov balancing problem.

\begin{theorem}[Matrix Lyapunov Calibration]
  \label{thm:matrix_lyapunov}
  Let $E$ be an $m$-dimensional target subspace. Suppose the regularized drift and diffusion on $E$ are given by $A_{\lambda,E} = -\lambda \Lambda_E$ and $D_{\lambda,E} = \frac{1}{\lambda} \Delta_E$, where $\Lambda_E, \Delta_E \in \mathbb{R}^{m \times m}$ are strictly positive definite matrices. Let $\Sigma_{\lambda,EE}$ be the unique positive definite solution to the matrix Lyapunov equation:
  \begin{equation}
    \label{eq:discussion_matrix_lyapunov}
    A_{\lambda,E}\Sigma_{\lambda,EE}
    +\Sigma_{\lambda,EE}A_{\lambda,E}^{*}
    +D_{\lambda,E}=0.
  \end{equation}
  Then $\Sigma_{\lambda,EE}$ scales exactly as $\lambda^{-2}$. Specifically, the normalized matrix $\widetilde{\Sigma}_{EE} \coloneqq \lambda^2 \Sigma_{\lambda,EE}$ is independent of $\lambda$ and is the unique positive definite solution to the scale-free Lyapunov equation:
  \begin{equation}
    \label{eq:scale_free_lyapunov}
    \Lambda_E \widetilde{\Sigma}_{EE} + \widetilde{\Sigma}_{EE} \Lambda_E^* = \Delta_E.
  \end{equation}
\end{theorem}
\begin{proof}
  Substituting $A_{\lambda,E} = -\lambda \Lambda_E$ and $D_{\lambda,E} = \frac{1}{\lambda} \Delta_E$ into \eqref{eq:discussion_matrix_lyapunov} yields:
  \[
    -\lambda \Lambda_E \Sigma_{\lambda,EE} - \lambda \Sigma_{\lambda,EE} \Lambda_E^* + \frac{1}{\lambda} \Delta_E = 0.
  \]
  Multiplying the entire equation by $\lambda$ gives:
  \[
    \Lambda_E (\lambda^2 \Sigma_{\lambda,EE}) + (\lambda^2 \Sigma_{\lambda,EE}) \Lambda_E^* = \Delta_E.
  \]
  Defining $\widetilde{\Sigma}_{EE} \coloneqq \lambda^2 \Sigma_{\lambda,EE}$, we recover \eqref{eq:scale_free_lyapunov}. Since $\Lambda_E$ is positive definite, its spectrum lies in the open right half-plane, making $-\Lambda_E$ Hurwitz. Because $\Delta_E$ is positive definite, standard continuous Lyapunov theory guarantees that \eqref{eq:scale_free_lyapunov} admits a unique positive definite solution $\widetilde{\Sigma}_{EE}$. Therefore, the exact scaling $\Sigma_{\lambda,EE} = \lambda^{-2} \widetilde{\Sigma}_{EE}$ holds for all $\lambda > 0$, properly generalizing the scalar calibration condition $\Sigma_{\lambda,kk} = \varepsilon_0 / (2\lambda^2)$.
\end{proof}

\subsection{Inverse-Free Residuals and Loewner Monotonicity}
\label{sec:discussion:operator-valued-residuals}

For an $m$-dimensional target space $E$, the Schur residual is a
positive-semidefinite operator on $E$. Since the free-block covariance
$C_\lambda \coloneqq \Sigma_{\lambda,II}$ acts as a compact, trace-class operator
on the infinite-dimensional space, it admits no bounded global inverse. Thus
the formal expression
\[
  \Sigma_{\lambda,EI}C_\lambda^{-1}\Sigma_{\lambda,IE}
\]
is not used as a definition.  The residual is defined instead by the
Anderson--Trapp variational shorting formula.

\begin{theorem}[Inverse-free operator residual and Galerkin Loewner monotonicity]
\label{thm:operator_residual}
Let $\mathcal H=E\oplus\mathcal H_I$, where $E$ is finite dimensional.
Let $C_\lambda=\Sigma_{\lambda,II}$ be a non-negative trace-class
covariance operator on $\mathcal H_I$, and set
\[
  B_\lambda\coloneqq \Sigma_{\lambda,IE}:E\to\mathcal H_I .
\]
Assume the Cameron--Martin range condition
\[
  B_\lambda E\subset H_{C_\lambda}
  \coloneqq \operatorname{Ran}(C_\lambda^{1/2}).
\]
Define the shorted residual $R_{\lambda,E}$ on $E$ by
\begin{equation}
\label{eq:discussion_operator_residual_form}
  \langle v,R_{\lambda,E}v\rangle_E
  =
  \sup_{\substack{u\in\mathcal H_I\\ \langle u,C_\lambda u\rangle>0}}
  \frac{|\langle u,B_\lambda v\rangle_{\mathcal H_I}|^2}
       {\langle u,C_\lambda u\rangle_{\mathcal H_I}},
  \qquad v\in E .
\end{equation}
Then $R_{\lambda,E}$ is a well-defined bounded positive-semidefinite
operator on $E$.  If
\[
  \|B_\lambda v\|_{H_{C_\lambda}}
  \le M\lambda^{-3}\|v\|_E,
  \qquad v\in E,
\]
then
\[
  \|R_{\lambda,E}\|_{\mathcal L(E)}=O(\lambda^{-6}).
\]

Let $(\mathcal H_n)_{n\ge1}$ be an increasing sequence of
finite-dimensional subspaces of $\mathcal H_I$ with dense union.  Define
the Galerkin residuals by the restricted variational problems
\begin{equation}
\label{eq:galerkin_variational_residual}
  \langle v,R_{\lambda,E}^{(n)}v\rangle_E
  \coloneqq
  \sup_{\substack{u\in\mathcal H_n\\ \langle u,C_\lambda u\rangle>0}}
  \frac{|\langle u,B_\lambda v\rangle_{\mathcal H_I}|^2}
       {\langle u,C_\lambda u\rangle_{\mathcal H_I}},
  \qquad v\in E .
\end{equation}
Then
\[
  0\preceq R_{\lambda,E}^{(n)}
  \preceq R_{\lambda,E}^{(n+1)}
  \preceq R_{\lambda,E}
\]
in the Loewner order, and
\[
  R_{\lambda,E}^{(n)}
  \longrightarrow
  R_{\lambda,E}
\]
in operator norm on $E$.
\end{theorem}

\begin{proof}
Fix $v\in E$ and write $y_v=B_\lambda v$.  By the range assumption,
$y_v\in H_{C_\lambda}$.  Thus the Cameron--Martin norm
$\|y_v\|_{H_{C_\lambda}}$ is finite.  Equivalently, the quadratic-form
Cauchy--Schwarz inequality gives
\[
  |\langle u,y_v\rangle_{\mathcal H_I}|^2
  \le
  \|y_v\|_{H_{C_\lambda}}^2
  \langle u,C_\lambda u\rangle_{\mathcal H_I},
  \qquad u\in\mathcal H_I .
\]
Hence the supremum in
\eqref{eq:discussion_operator_residual_form} is finite and equals
\[
  \|y_v\|_{H_{C_\lambda}}^2 .
\]
Therefore $v\mapsto \langle v,R_{\lambda,E}v\rangle_E$ defines a finite
non-negative quadratic form on the finite-dimensional space $E$, and is thus
represented by a unique positive-semidefinite operator $R_{\lambda,E}$.
The same argument applies to each restricted supremum and defines
$R_{\lambda,E}^{(n)}$.

Since
\[
  \mathcal H_n\subset \mathcal H_{n+1}\subset\mathcal H_I,
\]
the admissible class in the restricted variational problem increases
with $n$.  Hence, for every $v\in E$,
\[
  \langle v,R_{\lambda,E}^{(n)}v\rangle_E
  \le
  \langle v,R_{\lambda,E}^{(n+1)}v\rangle_E
  \le
  \langle v,R_{\lambda,E}v\rangle_E .
\]
This is precisely the Loewner monotonicity.

It remains to identify the limit.  Let
\[
  q_n(v)=\langle v,R_{\lambda,E}^{(n)}v\rangle_E,
  \qquad
  q(v)=\langle v,R_{\lambda,E}v\rangle_E .
\]
The sequence $q_n(v)$ is increasing and bounded above by $q(v)$.
Conversely, fix any admissible $u\in\mathcal H_I$ with
$\langle u,C_\lambda u\rangle>0$.  Since $\bigcup_n\mathcal H_n$ is
dense, there exist $u_n\in\mathcal H_n$ with $u_n\to u$ in
$\mathcal H_I$.  Because $C_\lambda$ is bounded and $y_v$ is fixed,
\[
  \langle u_n,C_\lambda u_n\rangle
  \to
  \langle u,C_\lambda u\rangle,
  \qquad
  \langle u_n,y_v\rangle
  \to
  \langle u,y_v\rangle .
\]
Thus every value of the full Rayleigh quotient can be approximated by
values taken over Galerkin subspaces.  Taking suprema gives
\[
  \lim_{n\to\infty}q_n(v)=q(v),
  \qquad v\in E .
\]
Since $E$ is finite dimensional, pointwise convergence of the quadratic
forms implies convergence of their representing matrices in operator
norm.

Finally, the assumed Cameron--Martin energy estimate gives
\[
  \langle v,R_{\lambda,E}v\rangle_E
  =
  \|B_\lambda v\|_{H_{C_\lambda}}^2
  \le
  M^2\lambda^{-6}\|v\|_E^2,
\]
and hence
\[
  \|R_{\lambda,E}\|_{\mathcal L(E)}
  \le M^2\lambda^{-6}.
\]
\end{proof}

\begin{remark}[Why strict Loewner growth is not expected]
\label{rem:strict_monotonicity}
The monotonicity in \Cref{thm:operator_residual} is non-decreasing
Loewner monotonicity.  Strict Loewner growth is not structurally
guaranteed.

Indeed, suppose that the Galerkin space is enlarged by a finite-dimensional
increment
\[
  \Delta\mathcal H_n
  =
  \mathcal H_{n+1}\ominus\mathcal H_n
\]
of dimension $k$.  In the Cameron--Martin energy geometry, the additional
Rayleigh--Ritz information available at the next step has rank at most
$k$.  Consequently
\[
  \operatorname{rank}
  \bigl(R_{\lambda,E}^{(n+1)}-R_{\lambda,E}^{(n)}\bigr)
  \le k .
\]
If the Galerkin scheme adds only one new direction at each step, then the
increment is at most rank one.  For a target space of dimension $m>1$, a
rank-one positive-semidefinite increment cannot be positive definite on
$E$.  Therefore
\[
  R_{\lambda,E}^{(n+1)}-R_{\lambda,E}^{(n)}\succ0
\]
is impossible for one-dimensional Galerkin updates when $m>1$.

Strict positivity of an increment requires both enough new Galerkin
directions, $k\ge m$, and full-rank visibility of the unresolved
cross-covariance defect across all target directions.
\end{remark}

\subsection{General observation operators and misalignment}
\label{sec:discussion:misalignment}

The scalar proof in this paper uses a block decomposition aligned with a single observed coordinate. However, typical finite-rank observation operators $L:\mathcal H\to\mathbb R^m$ select a subspace $E$ that does not commute with the free prior covariance, meaning the target-free split is not an eigenbasis split. This spectral misalignment generates cross-coupling that requires a precise resolvent analysis.

\begin{theorem}[Misaligned Covariance Coupling]
  \label{thm:misaligned_coupling}
  Let $P_E$ be an orthogonal projection onto an $m$-dimensional observed subspace $E$, and let $\mathcal{H}_I = E^\perp$. Let the regularized operator family be decomposed such that the target block matches \Cref{thm:matrix_lyapunov} with $A_{\lambda,E} = -\lambda \Lambda_E$, the cross-noise is zero ($D_{\lambda,IE} = 0$), and the off-diagonal drift coupling is denoted by $A_{\lambda,IE}$.
  Assume that, for some $\lambda_0>0$,
  \[
    \sup_{\lambda\ge\lambda_0}
    \|A_{\lambda,IE}\|_{\mathcal L(E,H_{Q_0})}<\infty.
  \]
  Assume also that the restricted drift $A_{II}$ generates an exponentially stable semigroup on $H_{Q_0}$ and that there exist $K,\gamma>0$ such that
  \[
    \|e^{-t\Lambda_E^*}\|_{\mathcal L(E)}
    \le
    K e^{-\gamma t},
    \qquad t\ge0.
  \]
  Then the cross-covariance satisfies the exact identity:
  \begin{equation}
    \label{eq:misaligned_resolvent_identity}
    \Sigma_{\lambda,IE} = \int_0^\infty e^{t A_{II}} A_{\lambda,IE} \Sigma_{\lambda,EE} e^{t A_{\lambda,E}^*} dt.
  \end{equation}
  The misaligned cross-covariance energy scales as $\|\Sigma_{\lambda,IE} v\|_{H_{Q_0}} = O(\lambda^{-3})$ for any $v \in E$.
\end{theorem}
\begin{proof}
  The off-diagonal block of the global Lyapunov equation $\mathcal{A}^{(\lambda)} \Sigma_\lambda + \Sigma_\lambda (\mathcal{A}^{(\lambda)})^* + \mathcal{D}^{(\lambda)} = 0$ yields the Sylvester equation:
  \[
    A_{II} \Sigma_{\lambda,IE} + \Sigma_{\lambda,IE} A_{\lambda,E}^* + A_{\lambda,IE} \Sigma_{\lambda,EE} = 0.
  \]
  Since $A_{II}$ generates an exponentially stable semigroup $e^{t A_{II}}$ and the assumed bound controls $e^{tA_{\lambda,E}^*}=e^{-\lambda t\Lambda_E^*}$, the unique solution is given by the integral representation \eqref{eq:misaligned_resolvent_identity}.

  To bound the Cameron--Martin norm, fix $v \in E$. We have $e^{t A_{\lambda,E}^*} v = e^{-\lambda t \Lambda_E^*} v$. By \Cref{thm:matrix_lyapunov}, $\Sigma_{\lambda,EE} = O(\lambda^{-2})$. By the uniform coupling assumption, choose $C$ such that $\|A_{\lambda,IE} w\|_{H_{Q_0}} \le C \|w\|_E$ for all $w \in E$ and all $\lambda\ge\lambda_0$. Using the uniform exponential stability of $e^{t A_{II}}$ on $H_{Q_0}$, namely $\|e^{t A_{II}}\|_{\mathcal{L}(H_{Q_0})} \le M e^{-\omega t}$ for some $M \ge 1, \omega > 0$, we estimate:
  \begin{align*}
    \|\Sigma_{\lambda,IE} v\|_{H_{Q_0}}
    &\le \int_0^\infty \|e^{t A_{II}}\|_{\mathcal{L}(H_{Q_0})} \|A_{\lambda,IE} \Sigma_{\lambda,EE} e^{-\lambda t \Lambda_E^*} v\|_{H_{Q_0}} dt \\
    &\le \int_0^\infty M e^{-\omega t} C \|\Sigma_{\lambda,EE}\|_{\mathcal{L}(E)} \|e^{-\lambda t \Lambda_E^*}\|_{\mathcal{L}(E)} \|v\|_E dt.
  \end{align*}
  Using $\|e^{-\lambda t\Lambda_E^*}\|_{\mathcal{L}(E)} \le K e^{-\lambda\gamma t}$, we obtain:
  \begin{align*}
    \|\Sigma_{\lambda,IE} v\|_{H_{Q_0}}
    &\le M C K \|\Sigma_{\lambda,EE}\|_{\mathcal{L}(E)} \|v\|_E \int_0^\infty e^{-(\omega + \lambda \gamma) t} dt \\
    &= M C K \|\Sigma_{\lambda,EE}\|_{\mathcal{L}(E)} \|v\|_E \frac{1}{\omega + \lambda \gamma}.
  \end{align*}
  Since $\|\Sigma_{\lambda,EE}\|_{\mathcal{L}(E)} = O(\lambda^{-2})$ and $(\omega + \lambda \gamma)^{-1} = O(\lambda^{-1})$, we conclude that $\|\Sigma_{\lambda,IE} v\|_{H_{Q_0}} = O(\lambda^{-3})$. Combined with \Cref{thm:operator_residual}, this verifies that the shorted-operator residual maintains the $O(\lambda^{-6})$ scaling even in the presence of observation misalignment.
\end{proof}

\subsection{Multidimensional Fredholm Determinant-and-Trace Cancellation}
\label{sec:discussion:multidim_functional}

With the finite-rank covariance bounds in
\Cref{thm:matrix_lyapunov,thm:operator_residual,thm:misaligned_coupling} in
place, the rank-one finite-part functional from \Cref{sec:gaussian_zero}
becomes a matrix problem.  The determinant contribution is no longer a scalar
logarithm expansion; it is a Fredholm determinant reduced to an
$m\times m$ matrix by the Weinstein--Aronszajn identity.  The residual trace terms
then cancels exactly against the matrix-valued shorted residual term.

The scalar regularized functional can be fully generalized to an $m$-dimensional observation subspace $E$. Let the exact evidence be $X_E = x_E^* \in \mathbb{R}^m$. The true Schur residual matrix is $S_{\lambda,E} \coloneqq \Sigma_{\lambda,EE} - R_{\lambda,E} \in \mathbb{R}^{m \times m}$, where $R_{\lambda,E}$ is the shorted-operator residual from \Cref{thm:operator_residual}. Since $R_{\lambda,E}^{(n)}\nearrow R_{\lambda,E}$ in the Loewner order, the Galerkin residual matrices
\[
S_{\lambda,E}^{(n)}
=
\Sigma_{\lambda,EE}-R_{\lambda,E}^{(n)}
\]
decrease in the Loewner order to
\[
S_{\lambda,E}
=
\Sigma_{\lambda,EE}-R_{\lambda,E}.
\]
They are not claimed to decrease strictly.  Positivity of the limiting
matrix is guaranteed by the condition
\[
\rho_{\lambda,E}
=
\Sigma_{\lambda,EE}^{-1/2}
R_{\lambda,E}
\Sigma_{\lambda,EE}^{-1/2}
<I.
\]

\begin{definition}[Multidimensional Fredholm Finite-Part Functional]
  \label{def:multidim_unified_action}
  Whenever $S_{\lambda,E} \succ 0$ and $\mu_{\text{obs},I}$ is equivalent to $\mu_{\lambda,I}$, the multidimensional Gaussian Fredholm finite-part functional is defined by evaluating the conditional residual in the metric of the shorted Schur matrix:
  \begin{equation}
    \label{eq:multidim_unified_action}
    \begin{aligned}
      \mathfrak{J}_{\lambda,E}(\mu_{\mathrm{obs}}; \mu_\lambda) \coloneqq &\KL(\mu_{\mathrm{obs},I} \parallel \mu_{\lambda,I}) \\
      &+ \frac{1}{2} \E_{\mu_{\mathrm{obs}}}\left[ \left( X_E - \mu_{E|I}(X_I) \right)^T S_{\lambda,E}^{-1} \left( X_E - \mu_{E|I}(X_I) \right) \right],
    \end{aligned}
  \end{equation}
  where $\mu_{E|I}(X_I) \coloneqq m_{\lambda,E} + W_{\Sigma_{\lambda,IE}}(X_I - m_{\lambda,I})$ is the measurable Cameron--Martin energy projection on $\mathbb{R}^m$, with the Wiener integral evaluated componentwise.
\end{definition}

To establish the multidimensional limit, we define the dimensionless coupling matrix $\rho_{\lambda,E} \coloneqq \Sigma_{\lambda,EE}^{-1/2} R_{\lambda,E} \Sigma_{\lambda,EE}^{-1/2} \in \mathbb{R}^{m \times m}$. By the energy bound in \Cref{thm:operator_residual}, $\|\rho_{\lambda,E}\| = O(\lambda^{-4})$.

\begin{theorem}[Exact Fredholm Trace and Matrix Determinant Cancellation]
  \label{thm:multidim_cancellation}
  For any evidence level $t \in \mathbb{R}^m$, let $\mu_{\lambda}^t$ be the weakly continuous conditional law on $X_E = t$. The infinite-dimensional KL divergence and the matrix-valued shorted residual expectation undergo exact cancellation of all cross-coupling traces and quadratic forms, yielding:
  \begin{equation}
    \label{eq:multidim_exact_cancellation}
    \mathfrak{J}_{\lambda,E}(\mu_\lambda^t; \mu_\lambda) = \frac{1}{2} (t - m_{\lambda,E})^T \Sigma_{\lambda,EE}^{-1} (t - m_{\lambda,E}) - \frac{1}{2}\log\det(I-\rho_{\lambda,E}).
  \end{equation}
\end{theorem}
\begin{proof}
  Let $C_\lambda \coloneqq \Sigma_{\lambda,II}$. Under $\mu_\lambda^t$, the free coordinates follow $\mathcal{N}(\mu_{I|E}(t), \Sigma_{I|E})$. By the Feldman--H\'{a}jek formula for trace-class perturbations, evaluating the KL divergence relies on the operator $K \coloneqq C_\lambda^{-1/2} \Sigma_{\lambda,IE} \Sigma_{\lambda,EE}^{-1/2}$. By the Weinstein--Aronszajn identity, the non-zero spectrum of the infinite-dimensional operator $KK^*$ matches that of the $m \times m$ matrix $K^*K = \rho_{\lambda,E}$. This reduces the infinite-dimensional Fredholm determinant and trace exactly to the matrix manifold:
  \begin{equation*}
    \begin{aligned}
      \KL(\mu_{\lambda,I}^t \parallel \mu_{\lambda,I}) = \frac{1}{2} \bigg[ &(t - m_{\lambda,E})^T \Sigma_{\lambda,EE}^{-1} R_{\lambda,E} \Sigma_{\lambda,EE}^{-1} (t - m_{\lambda,E}) \\
      &- \log\det(I - \rho_{\lambda,E}) - \tr(\rho_{\lambda,E}) \bigg].
    \end{aligned}
  \end{equation*}
  For the residual expectation term, let $Z \coloneqq t - \mu_{E|I}(X_I)$. Using block matrix identities, the conditional mean is $\E_{\mu_\lambda^t}[Z] = S_{\lambda,E} \Sigma_{\lambda,EE}^{-1} (t - m_{\lambda,E})$ and the conditional variance is $\Var_{\mu_\lambda^t}(Z) = R_{\lambda,E} \Sigma_{\lambda,EE}^{-1} S_{\lambda,E}$. Expanding the quadratic form $\frac{1}{2}\E[Z^T S_{\lambda,E}^{-1} Z]$ gives:
  \begin{equation*}
    \begin{aligned}
      \frac{1}{2} \bigg[ &(t - m_{\lambda,E})^T \Sigma_{\lambda,EE}^{-1} (t - m_{\lambda,E}) \\
      &- (t - m_{\lambda,E})^T \Sigma_{\lambda,EE}^{-1} R_{\lambda,E} \Sigma_{\lambda,EE}^{-1} (t - m_{\lambda,E}) + \tr(\rho_{\lambda,E}) \bigg].
    \end{aligned}
  \end{equation*}
  Summing the KL divergence and the quadratic residual form exactly cancels both the $\pm \frac{1}{2}\tr(\rho_{\lambda,E})$ residual trace terms and the cross-covariance residual matrices, isolating \eqref{eq:multidim_exact_cancellation}.
\end{proof}

\begin{theorem}[Multidimensional Fredholm Finite-Part Regularization]
  \label{thm:multidim_regularization}
  Assume the $m$-dimensional regularized drift on $E$ is centered at the observed value:
  \[
    A_{\lambda,E}(m_{\lambda,E}-x_E^*)+b_E=0,
    \qquad
    A_{\lambda,E}=-\lambda\Lambda_E.
  \]
  Defining the normalized displacement $\delta_E \coloneqq \Lambda_E^{-1} b_E \in \mathbb{R}^m$, the mean deterministic offset satisfies
  \[
    m_{\lambda,E}-x_E^*
    =
    \lambda^{-1}\Lambda_E^{-1}b_E
    =
    \lambda^{-1}\delta_E.
  \]

  As $\lambda \to \infty$, the multidimensional Fredholm finite-part functional converges to a well-defined bounded limit governed by the scale-free Lyapunov solution $\widetilde{\Sigma}_{EE}$:
  \begin{equation}
    \lim_{\lambda\to\infty} \mathfrak{J}_{\lambda,E}(\mu_{\mathrm{obs}}; \mu_\lambda) = \frac{1}{2} \delta_E^T \widetilde{\Sigma}_{EE}^{-1} \delta_E.
  \end{equation}
  The approximation residual satisfies $\mathfrak{J}_{\lambda,E} = \frac{1}{2} \delta_E^T \widetilde{\Sigma}_{EE}^{-1} \delta_E + O(\lambda^{-4})$.
\end{theorem}
\begin{proof}
  Substitute $t = x_E^*$ and $t - m_{\lambda,E} = -\lambda^{-1} \delta_E$ into \Cref{thm:multidim_cancellation}. Under the canonical scaling,
  \[
    \Sigma_{\lambda,EE}
    =
    \lambda^{-2}\widetilde\Sigma_{EE},
    \qquad
    t-m_{\lambda,E}
    =
    -\lambda^{-1}\delta_E,
    \qquad
    \|\rho_{\lambda,E}\|=O(\lambda^{-4}).
  \]
  Therefore the quadratic term yields exactly
  \[
    \frac12
    \delta_E^T
    \widetilde\Sigma_{EE}^{-1}
    \delta_E,
  \]
  and the determinant term satisfies $\log\det(I-\rho_{\lambda,E})=O(\lambda^{-4})$. Hence
  \[
    \mathfrak{J}_{\lambda,E}
    =
    \frac12
    \delta_E^T
    \widetilde\Sigma_{EE}^{-1}
    \delta_E
    +
    O(\lambda^{-4}).
  \]
\end{proof}

\begin{remark}[Algebraic Mechanism of Uniform Positivity]
  \label{rem:asymptotic_coupling_precision}
  The $O(\lambda^{-4})$ decay of the normalized coupling matrix $\|\rho_{\lambda,E}\|$ is an asymptotic bound providing the algebraic mechanism that ensures the uniform positivity of the Fredholm finite part. By extracting the leading-order expansion of the cross-covariance resolvent under the canonical scaling $D_{\lambda,E} = \Delta_E / \lambda$, the trace of the coupling matrix converges exactly as:
  \begin{equation}
    \lim_{\lambda\to\infty} \lambda^4 \tr(\rho_{\lambda,E}) = \frac{\omega}{d_{\mathrm{bg}}} \sum_{j=1}^m \frac{c_j^2 \Delta_j}{\Lambda_j^3}.
  \end{equation}
  
  This extreme asymptotic decoupling directly resolves the algebraic obstruction of non-strict Loewner growth noted in \Cref{rem:strict_monotonicity}. The Anderson--Trapp formulation and its Galerkin approximations are employed purely to establish the existence of the limiting shorted operator on the Hilbert space via the Monotone Sequence Theorem. For any nested subspace sequence, the finite-dimensional operator residuals satisfy $0 \preceq R_{\lambda,E}^{(n)} \preceq R_{\lambda,E}$. Normalizing by the target covariance block $\Sigma_{\lambda,EE} = O(\lambda^{-2})$ yields the spectral bound:
  \[
    I \succeq \Sigma_{\lambda,EE}^{-1/2} S_{\lambda,E}^{(n)} \Sigma_{\lambda,EE}^{-1/2} \succeq I - \rho_{\lambda,E}.
  \]
  Because the limiting normalized coupling strictly vanishes as $\tr(\rho_{\lambda,E}) = O(\lambda^{-4})$, the spectrum of every intermediate operator $S_{\lambda,E}^{(n)}$ is uniformly bounded from below by $1 - O(\lambda^{-4}) > 0$. The sequence of conditional covariance operators is therefore uniformly positive definite. This guarantees that the Gaussian conditioning remains well-posed along the entire approximation path, demonstrating that shorted operators provide the necessary and sufficient algebraic mechanism to define the Fredholm finite part at singular evidence, without requiring strict Loewner monotonicity. The exact $O(\lambda^{-4})$ scaling of the strict coupling trace is verified numerically to arbitrary precision in \Cref{app:experiments}.  The necessity of extended-precision arithmetic is a direct consequence of the spectral ill-conditioning (Theorem~\ref{thm:triple_impossibility}\ref{item:numerical_impossibility}): at $N=2048$, $\kappa\sim4\times10^{6}$ destroys double-precision accuracy, and the pre-asymptotic spectral regime forces $\lambda$ below the threshold where $\lambda^{4}\tr(\rho_{\lambda,E})$ is resolvable in IEEE-754 double precision.  Quadruple-precision (256-bit) arithmetic is therefore required to extract the true $O(\lambda^{-4})$ decay.
\end{remark}

\subsection{Non-Gaussian and time-dependent extensions}
\label{sec:discussion:non-gaussian}

The result is deliberately limited to finite-rank covariance block algebra and
its Gaussian Radon application.  It does not assert a general non-Gaussian
conditioning theory.  For non-Gaussian stationary processes, the projection
residual $S_{\lambda,E}=\Sigma_{\lambda,EE}-R_{\lambda,E}$ still has a
second-moment interpretation, but Radon disintegration, finite-part
functionals, and measure-class behavior require separate hypotheses
\cite{bogachev2007measure,kallenberg2021foundations}.  Time-dependent relaxation
after conditioning, or transport distances between post-conditioning laws,
belongs to a different problem.  The contribution here is the
operator-theoretic removal of the compact-block Schur obstruction.
